\documentclass[../Lael_John_MasterThesis_Draft.tex]{subfiles}
\graphicspath{\subfix[../images]}
\begin{document}
Since the statement of the conjecture in 2008, there has been a substantial amount of work done to arrive at a conclusive final resolution to the problem. Some of this body of work relies on the specific geometric construction of secant varieties and how they relate to the question of computing rank and symmetric rank. This connection to geometry is not surprising because tensors can also be seen as points lying in some topological space (much like vectors and matrices can be seen as representing points, as well as algebraic objects.) This chapter goes over some of the work on proving and disproving Comon, with the final section of this chapter compiling all the conditions currently known for which the conjecture holds or doesn't hold.

% \section{Secant Varieties}\label{section:secantvarieties}
% The notion of studying secants to a given geometric object is not novel, but doing so sheds more light on the description of tensors as lying in projective space. Considering a curve in a vector space, the secant line joining any two distinct points is described exactly by all the linear combinations of the points. This analogy is the inspiration for the following definition.
% \begin{figure}
%     \centering
%     \begin{tikzpicture}[scale=0.3]
%         % Coordinates
%         \coordinate (A) at (4.51, 6.52);
%         \coordinate (B) at (4.51, 6.52);
%         \coordinate (C) at (1.00, 2.50);
%         \coordinate (D) at (9.25, 11.00);
%         \coordinate (E) at (2.75, 0.75);
%         \coordinate (F) at (8.75, 12.00);
%         \coordinate (G) at (3.75, 0.25);
%         \coordinate (H) at (10.50, 9.00);

%         % Drawing
%         \draw (A) circle (3.25);
%         \draw (C) -- (D);
%         \draw (E) -- (F);
%         \draw (G) -- (H);
%     \end{tikzpicture}
%     \caption{Secant lines through a circle (and a limiting tangent line)}\label{figure:secantexample}
% \end{figure}
% \begin{definition}[$r$-th Secant Variety to $X$]\label{defn:secantvariety}
%     Let $X \subset \proj V$ be a projective variety. Then the $r$-th secant variety to $X$, denoted $\sigma_r(X)$ is defined as $$\sigma_r(X) := \overline{\bigcup_{x_1, \ldots, x_r \in X} \proj_{x_1, \ldots, x_r}}$$Here $\proj_{x_1, \ldots, x_r}$ (for $x_1, \ldots, x_r$ lying in general position) is the projectivisation of the $r$-dimensional hyperplane defined by these $r$ points. The bar denotes the Zariski closure in $\proj V$, covering the cases where the $r$ points chosen may not lie in general position, particularly also when all the $r$ chosen points coincide leading to the corresponding tangent hyperplane.
% \end{definition}
% \begin{observation}
%     It is clear from the definition above that successive secant varieties are nested, with $\sigma_0(X) = X$, $\sigma_1(X)$ describing all the tangent lines to $X$, and so on, such that $\sigma_{r-1}(X) \subseteq \sigma_r(X)$. Such a sequence of nested geometric objects must stabilise, for reasons of dimension. 
% \end{observation}
% \begin{remark}
%     While this geometric object is easy to describe with the notation above, computing the ideals for such a variety, or even the equations required to cut such a variety out form about half of the body of work surrounding Comon's conjecture. The following papers \cite{buczynska_secant_2013,buczynski_ranks_2013,bernardi_waring_2020} develop alternative geometric tools to secant varieties that are easier to compute.
% \end{remark}
% Since elementary tensors can be seen as lying on either the Segre or Veronese embedding, the $r$-th secant varieties to each of these describe all the tensors of \textit{border rank} at most $r$. Note here that it is border rank that because of the closure in the definition of the secant variety. To describe all tensors of \textit{rank} at most $r$, it is enough to look at the \textit{interior} points of the $r$-th secant variety, this interior again being under the Zariski topology on $\proj V$. 

% % In fact \cite{bernardi_hitchhiker_2018} provides a comprehensive guide to the study of secant varieties, and highlights the importance of the Alexander-Hirschowitz theorem in computing the dimension of secant varieties. As previously observed, since secant varieties are nested, finding the first secant variety that covers all points in projective space allows for an upper bound on the length of all decompositions of tensors. This allows for the computation of "generic" rank

\section{Attempts to Prove the Conjecture}\label{section:pastoverview}
The statement of Comon's Conjecture can now be stated in more formal terminology, as introduced in the previous chapter.

\begin{conjecture}[Comon et al.\cite{comon_symmetric_2008}]
    Given a symmetric, order-$d$ tensor $T \in S^d(\field^n) \subset (\field^n)^{\otimes d}$, then $\rank(T) = \srank(T)$.
\end{conjecture}

The ability to rewrite symmetric tensors of fixed order both as homogenous polynomials and as a sum of regular tensors was the key factor that seems to have inspired this conjecture, as stated in subsection 4.4 from Comon et al. \cite{comon_symmetric_2008}. This section details some of the papers that prove that the conjecture holds under certain constraints.

As mentioned in the introduction Comon et al. \cite{comon_symmetric_2008} showed that the conjecture was true for tensors of "large enough" order --- a notion that became concrete in later literature.
[COMON ET AL. 2008 PROPOSITION 5.4 - PROOF OF SUFFICIENTLY LARGE (include proposition 4.3 here) CAUSING REWRITE OF THE FOLLOWING PARAGRAPH]
To establish such a proof, the authors first proved the following lemma (stated as a proposition), highlighting how a finite collection of linear polynomials, under pairwise independence constraints could be used to generate a linearly independent set of degree-$d$ homogenous polynomials. Doing so requires the following definition:
\begin{definition}[Bilinear form on the space of homogenous polynomials \cite{comon_symmetric_2008}]
	Let $F, G \in \field[x_1, \ldots, x_n]_d$, where $$F(\mathbf{x}) = \sum_{|\alpha| = d} \binom{d}{a_1,\ldots, a_n}\lambda_{\alpha}\mathbf{x}^{\alpha},\quad G(\mathbf{x}) = \sum_{|\alpha| = d} \binom{d}{a_1,\ldots, a_n}\mu_{\alpha}\mathbf{x}^{\alpha}.$$Here $\alpha = (a_1, \ldots, a_n)$ and $\lambda_{\alpha}, \mu_{\alpha} \in \field$ for all $\alpha$. Then the form $\langle \cdot, \cdot \rangle: \field[x_1, \ldots, x_n]_d \times \field[x_1, \ldots, x_n]_d \to \field$ defined as $$\langle F, G \rangle := \sum_{|\alpha| = d} \binom{d}{a_1,\ldots,a_n}\lambda_{\alpha}\mu_{\alpha},$$is a symmetric, non-degenerate bilinear form. 
\end{definition}
\begin{lemma}[\cite{comon_symmetric_2008}]\label{lemma:linearformsindependence}
	Let $L_1, \ldots, L_r \in \field[x_1,\ldots, x_n]_1$, be linear forms in $n$ variables. Suppose for all $i \neq j$, $L_i \neq \alpha L_j$ (where $\alpha \in \field$). Then for any $d \geq r-1$, the polynomials $L_1^d, \ldots, L_r^d$ are linearly independent in $\field[x_1, \ldots, x_n]$. 
\end{lemma}
\begin{proof}
	Fix $d \geq r-1$. Suppose if possible, that there exist coefficients $\alpha_1, \ldots \alpha_r \in \field$ such that $\sum_{i=1}^{r} \alpha_iL_i^d = 0$.  Since the bilinear form is non-degenerate, the expression $$ \langle F, \sum_{i=1}^{r} \alpha_iL_i^d \rangle = 0$$for any $F \in \field[x_1, \ldots, x_n]_d$. This simplifies to $$ \sum_{i=1}^{r} \alpha_i\langle F, L_i^d \rangle = \sum_{i=1}^{r} \alpha_i F(L_i) = 0$$where the expression $F(L_i) = F(\beta_1, \ldots, \beta_n)$ when $L_i = \sum_{i=1}^{n} \beta_ix_i$. The next step is to show that a particular degree-$d$ homogenous polynomial $F$ exists, such that it vanishes on all except one of the $L_i$. Consider $F$, a degree-$d$ homogenous polynomial, to be the multiple of a product of $r-1$ linear forms, such that $F(L_r) \neq 0$ but $F(L_i) = 0$ for all $i < r$. Then one can determine that $\alpha_r = 0$, by the bilinear form expansion derived above. In a similar manner, similar manner, it can be shown that all the other $\alpha_i$ are also zero, leading to the conclusion that the forms $L_i^d$ are linearly independent.
\end{proof}
This lemma also allows for determining when symmetric, order-$d$ tensors can be considered as linearly independent, using the isomorphsim in \Cref{prop:homotensorisomorphism}.
\begin{proposition}[\cite{comon_symmetric_2008}]
	Let $v_1, \ldots, v_s \in \field^n$ be pairwise linearly independent. If $d$ is sufficiently large, then the symmetric tensor, defined as $$T := \sum_{i=1}^{s} v_i^{\otimes d}$$satisfies $\rank(T) = \srank(T)$ generically (i.e in a dense subset of order-$d$ symmetric tensors).
\end{proposition}
\begin{proof}
	To begin with, since $S^d(\complex^n) \subseteq (\complex^{n})^{\otimes d}$, for all $T \in S^d(\complex^n),  \rank(T) \leq \srank(T)$. It is enough to show that in this specific case, for a tensor $T$ as defined in the statement of this proposition, the reverse inequality i.e. $\rank(T) \geq \srank(T)$ holds.

	Suppose $\rank(T) = r$ and $\srank(T) = s$. Then there exist decompositions $$\sum_{i=1}^{r} x_1^i \otimes \cdots \otimes x_d^i = T = \sum_{j=1}^{s} v_j^{\otimes d}.$$ Without a loss of generality, assume that $d$ is even. Then $T$ can be written instead as $T = \sum_{i=1}^{s} v_j^{\otimes d/2} \otimes v_j^{\otimes d/2}$. By the previous lemma, for $d/2 \geq s-1$, the tensors $v_j^{\otimes d/2}$ are generically linearly independent. Now therefore, one can find functionals $\Phi_1, \ldots, \Phi_s \in S^{d/2}(\field^n)^*$ that are duals to the vectors $v_j^{\otimes d/2}$. i.e. $\Phi_i(v_j^{\otimes d/2}) = \delta_{i,j}$, the Kronecker delta. Then the initial equality of decompositions can be transformed (after contracting along the first $d/2$ modes) into the following expressions$$\sum_{i=1}^{r} \alpha_{ij}x_{d/2 + 1}^i\otimes \cdots \otimes x_{d}^i = v_j^{\otimes d/2},$$for each of the $v_j$, where $\alpha_{ij} = \Phi_j(x_1^i \otimes \cdots \otimes x_{d/2}^i)$ where $i = 1, \ldots, r$. This then implies that the $s$ linearly indepenedent vectors $v_j^{\otimes d/2}$ lie in the span of the $r$ tensors of the form $\{x_{d/2 + 1}^i \otimes \cdots \otimes x_d^i\}_{i=1}^r$. This then implies that $r \leq s$, and the proof is complete.
\end{proof}

The authors further noted that, as a result of the Veronese embedding (a canonical construction in algebraic geometry), the rank of a symmetric tensor $T \in S^d(\complex^n)$ cannot be more than $\binom{n + d - 1}{d}$. Another major result in their paper was that the conjecture held generically whenever $\srank(T) \leq n$, the dimension of the underlying vector space. In particular, they were able to show that for every $T\in S^d(\complex^n)$, the conjecture held if $\srank(T) = 1$ or $2$.

% Here the term generic (and typical) have specific mathematical meanings. Generic is used for a property that holds almost everywhere, whereas typical properties hold on non-zero volume sets. Such a distinction is not particularly pertinent over $\complex$, but may vary over other fields. 
% Brachat et al. \cite{brachat_symmetric_2010}, in 2010 proposed a new algorithm for symmetric tensor decomposition beginning with the relationship between symmetric tensor decomposition and the secant varieties of a Veronese variety (see Section \cref{section:secantvarieties}). Their resulting algorithm required a reformulation of the problem in a dual space, and Hankel operators.[INCLUDE THIS?]

A concurrent approach to tackling Comon's conjecture arises from the transposition of a classical problem from number theory --- the Waring decomposition of a number --- to the language of polynomials. This famous problem then gets restated as follows:
\vspace{12pt}
\hrule
\vspace{12pt}
\textit{"What is the the minimum number of linear polynomials required such that a linear combination of them to the $d$-th power is equal to a homogenous degree-$d$ polynomial?"}
\vspace{6pt}
\hrule
\vspace{12pt}
In mathematical notation, take $f \in \field[x_1, \ldots, x_n]_d$, a homogenous polynomial of degree $d$ and $\mathbf{a}:= (a_1, \ldots, a_n)$ denoting a vector containing coefficients such that $\mathbf{a}\cdot\mathbf{x} := a_1x_1 + \cdots + a_nx_n$. Then the Waring problem for polynomials is the following minimisation problem \begin{equation}
	f:= \sum_{|\alpha| = d} f_{\alpha}\mathbf{x}^{\alpha} = \min_{s > 0} \sum_{i=1}^s(\mathbf{a}_i\cdot \mathbf{x})^d
\end{equation} 

A solution to the polynomial Waring problem solves the issue of computing symmetric rank for a given symmetric tensor (see \Cref{subsection:homogenouspolyrelation}). In fact this relationship is so significant that in the literature, the term \textbf{Waring Rank} is used synonymously with symmetric rank when discussing symmetric tensors. An equivalent formulation for the Waring problem is in terms of predicting the dimensions of higher secant varieties and relating these predictions to the symmetric (or Waring) decompositions for symmetric tensors \cite{bernardi_hitchhiker_2018}. 

The polynomial Waring problem for generic forms was solved by Alexander and Hirschowitz in 1995, by using this secant variety formulation, building on previous work from Campbell, Palatini and Terracini, and using the novel "differential Horace's method" \cite{bernardi_hitchhiker_2018}. They showed that $\sigma_r(\nu_d(\proj \field^{n+1}))$ - the $r$-th secant variety to the $d$-th Veronese embedding of $\field^{n+1}$ - is of expected dimension with 5 exceptional cases. Brambilla and Ottaviani \cite{brambilla_alexanderhirschowitz_2008} in 2008 gave a self-contained proof of the theorem, with some simplifications, particularly in the case of order $3$ symmetric tensors. The existence of concrete methods for computing rank and symmetric rank meant that for any given symmetric tensor, Comon's conjecture could be quickly verified or disproved.

% Landsberg \cite{landsberg_tensors_2012}, in his textbook describes the conjecture in terms of (border) rank preserving pairs of projective varieties. 
% \begin{definition}[(Border) Rank Preserving Pair]\label{defn:brpp}
%     $X \subseteq \proj V$, a variety, and $L \subseteq \proj V$ a linear subspace when taken together as $(X,L)$ are a \textit{(border) rank preserving pair} if $\langle(X \cap L)_{red}\rangle = L$ and $\rank(X|_L) = \rank(X|_{(X \cap L)_{red}})$. (A similar condition holds for border rank, but this can now be stated in terms of secant varieties as $\sigma_r(X) \cap L = \sigma_r((X\cap L)_{red}))$. Here $(X\cap L)_{red}$ refers to the intersection without any points of non-unit multiplicity. and $\rank$ is now seen as a function from $\langle X \rangle \to \naturals$
% \end{definition}
% \begin{remark}
%     This translates the statement of Comons conjecture from algebra to a statement about (secant) varieties lying in some projective space. In this restatement, $\Sigma(\proj \field^n \times \cdots \times \proj \field^n)$ (the $d$-fold Segre embedding) and $\proj(S^d(\field^n))$ (all the lines passing through the $d$-th Veronese) are required to be a rank-preserving pair.
% \end{remark}
% One of the major issues with trying to work on Comon's conjecture via the world of geometry is the fact that computing the defining equations for secant varieties involves computin determinants or minors, which is computationally quite expensive/prohibitive.

Ballico and Bernardi \cite{ballico_tensor_2013} in 2013 showed that that a tensor $T$ lying on the tangent variety to a Veronese variety must always have $\rank(T) = \srank(T)$. Their paper highlighted explicit algorithms for the computation of the rank of a tensor that could be expressed as the limit of a linear combination of elementary tensors. The main theorem in their paper is as follows. The Segre variety as a classical construction in algebraic geometry is used to visualise elementary tensors that live in a tensor product of vector spaces. The Veronese variety (again a classical construction), can be seen as embedded into this larger variety, showing only the elementary symmetric tensors. Thus general points in the spaces containing the Segre (or Veronese) can be seen as general tensors. 
\begin{theorem}[\cite{ballico_tensor_2013}]
	Let $\tau(X_{n_1, \ldots, n_d})$ be the tangential variety of the Segre variety $X_{n_1, \ldots, n_d}$. Then, for each $P \in \tau(X_{n_1, \ldots, n_d})$, one has that $$\rank_{X_{n_1, \ldots, n_d}} = \eta_{X_{n_1, \ldots, n_d}}(P),$$ where $\eta_{X_{n_1, \ldots, n_d}}(P)$ denotes the minimum size of a set $E \subseteq \{1, \ldots, d\}$ such that $P$ as a linear subspace lies in the linear span of exactly $|E|$ linear subspaces, passing through a point $O \in X_{n_1, \ldots, n_d}$
\end{theorem}
\begin{remark}
	This theorem equates the minimum number of points lying anywhere on the Segre variety whose linear span contains the point $P$ lying in the tangential variety, and the minimum number of linear spaces generated at a point $O$ on the Segre in whose linear span $P$ must lie.
	The application to Comon's Conjecture comes from the following corollary
\end{remark}
\begin{corollary}
	Suppose $P$ lies in the linear span of any two points lying on the Veronese (embeddable inside the Segre). Then $\rank(P) = \srank(P)$.
\end{corollary}


% In the literature reviewed, one of the most popular techniques used to compute the symmetric rank of a given order $d$ tensor is the construction of \textit{secant varieties} to the Veronese variety $\nu_d(\proj \field^n)$ --- since this variety describes all the elementary order-$d$ symmetric tensors geometrically. Points lying on the $r$-th secant variety to such a construction are guaranteed to have symmetric rank $r$ (Exceptions are covered by the Alexander-Hirschowitz theorem, that describes the discrepancy between expected rank and actual rank for symmetric tensors \cite{brambilla_alexanderhirschowitz_2008}.) 

% IF NOT INCLUDING SECANT VARIETIES, WHY INCLUDE THIS? JUSTIFY.[While secant varieties are convenient geometric tools to use for computing rank and symmetric rank via the Alexander-Hirschowitz theorem, Buczynska and Buczynski in 2013 \cite{buczynska_secant_2013} demonstrated how the commonly used practice of using the $(r+1) \times (r+1)$ minors of the \textit{catalecticant matrix} to compute the generating equations for the secant variety are not sufficient for higher order Veronese varieties (order $>> 2$). This paper resulted in the developement of \textit{cactus varieties}.] 

% Oeding and Ottaviani \cite{oeding_eigenvectors_2013} in 2013 contributed new algorithms for computing solutions to the polynomial Waring decomposition problem - equivalent to the computation of secant varieties to a given Veronese variety. They did so by further developing work done to define the eigenvectors of a tensor. Using these eigenvectors and Kozul flattenings allowed for an algorithm that could compute the Waring decomposition while also remaining valid for previously covered cases under the catalecticant matrix methodology.

Computing decompositions covers the first question about symmetric rank --- the existence of a minimal length symmetric decomposition. The other question, concerning the uniqueness of a decomposition,  requires the definition of when a tensor is \textit{identifiabililty}. 

\begin{definition}[Identifiability]\label{defn:identifiability}
    A tensor $T \in V_1 \otimes \cdots \otimes V_d$ is said to be \textbf{identifiable} if it admits a unique rank decomposition (modulo scaling by $d$-th roots of unity and permutations of summands)
\end{definition}
\begin{example}
	The tensor $(3 , 4)\otimes (1, 1) + (4, 5) \otimes (1, 0)$ can be written as $(7, 9)\otimes (1, 0) + (3, 4) \otimes (0, 1)$, and is thus not identifiable, since it can be decomposed in distinct ways, not involving simply a permutation of summands or a direct rescaling.
\end{example}
% \begin{observation}
%     This definition is equivalent to asking that an identifiable tensor lie on only one secant variety (again modulo permutations of points on the Segre embedding, and scaling factors)
% \end{observation}

% When dealing with questions of tensor rank specifically, given an identifiable tensor, it suffices to find only one rank decomposition. This then becomes useful for the purposes of looking for a counterexample to Comon's conjecture. The likely candidate symmetric tensors must either have a symmetric rank greater than the generic rank, or must lie in the complement of a dense subset. [REFINE THIS SENTENCE] If Comon's conjecture does not hold

Generic identifiability is the idea that almost every tensor (i.e. a tensor in a dense open subset of tensors), is identifiable. Specific identifiability poses the opposite question - given a specific tensor $T$, does it admit a unique decomposition? Chiantini, Ottaviani and Vannieuwenhoven \cites{chiantini_generic_2016, chiantini_effective_2017} in two subsequent papers tackled both of these questions. Their earlier paper showed that a general symmetric tensor in $S^d(\complex^{n+1})$ of sub-generic rank $r$ was identifiable, provided $r \leq \binom{n+d}{d}/(n+1)$ (with all exceptions again being covered by the Alexander-Hirschowitz theorem). Their second paper focused on showing what criteria could be called "effective" to determine if a given tensor was specifically identifiable. Criteria were designated effective if their conditions were satisfied by generic points (i.e. the criteria were satisfied in an open, dense subset). One of the conclusions of this last paper was that Kruskal's criterion (a condition on collections of order-$d$ tensors, seen as points in a vector space) was \textit{effective}. Following Section 7 in \cite{chiantini_effective_2017}, the authors address the applications of their results to Comon's conjecture via the following theorem
\begin{theorem}[\cite{chiantini_effective_2017}]
	Let $$T = \sum_{i=1}^{r}\phi_iv_i^{\otimes d},$$ with $\phi_i \in \field\setminus\{0\}$, and $v_i \in \field^n$. Then $T$ is a generic, order-$d$, symmetric tensor, lying in $(\field^n)^{\otimes d}$. If \begin{equation*}
		r \leq \begin{cases}
			\frac{3}{2}n -1 &  d = 3\\
			\binom{k+n}{n} + \frac{1}{2}\binom{2+n}{n} -1 & d = 2k + 1\\
			\binom{k+n}{n} -n -1 & d = 2k,
		\end{cases}
	\end{equation*}
	with $2 \leq k \in \naturals$, then $T$ admits only symmetric (or Waring) decompositions, and the symmetric and tensor rank of $T$ must coincide
\end{theorem}

Friedland \cite{friedland_remarks_2016} in 2016 focused not on the computation of the symmetric rank, but on the properties of symmetric tensors of a given symmetric rank. In particular the paper focused on showing that when a symmetric tensor $T$ can be flattened into a matrix $T^{(i)}$, then Comon's Conjecture held if the rank of the tensor, $\rank(T) = \rank(T^{(i)})$ or if $\rank(T) = \rank(T^{(i)})$. [PROVIDE FULL STATEMENT FOR FRIEDLAND THEOREM 1.1. COMBINE SECTIONS 3,5,6.]
\begin{theorem}[\cite{friedland_remarks_2016}]
	Let $d \geq 3$, and $|\field| \geq 3$. Let $T \in S^d(\field^n)$. Suppose that $\rank(T) \leq \rank(A(T)) + 1$. Then $\srank(T) = \rank(T)$.
\end{theorem} 
\nomenclature[D]{$T^{(i)}$}{The flattening of a tensor $T$ along the $i$th mode (see \Cref{defn:tensorflattening})}
Zhang et al. \cite{zhang_comons_2016} conclusively showed that Comon's conjecture held for all tensors where the rank of the tensor under consideration is not more than its order. [INCLUDE THEOREM 3.8 AND PROOF]

% Carlini et al. \cite{carlini_real_2017} in 2017 discussed the computation of real and complex rank for symmetric tensors associated to monomials of a fixed degree, and showed that these must agree when the least exponent in the monomial is 1. They did so by using duals to the homogenous ring of polynomials to describe a differentiation operation, giving rise to ideals that were "perpendicular" to any one monomial under scrutiny.

[REMARK 3.11 AFTER COROLLARY 3.10] Ballico et al \cite{ballico_bounds_2018} at the end of 2018, refocused their efforts on the question of identifiability, and worked on refining Kruskals criterion (sharp and algebraic in nature) into a geometric criterion that can be verified more easily. They concluded that Comon's conjecture held for tensors whose symmetric rank $r$ is strictly bounded above by $\binom{n+e}{e}$ for a specific choice of $e$. This choice depended upon splitting the Segre product (another classical construction from algebraic geometry) into two disjoint components and then imposing conditions on a set of points on the Segre variety that exhibited distinct behaviour when projected down to the first component versus the second. The authors of this paper clearly pointed out that even though they had increased the scope of validity of Comon's conjecture, Shitov's proposed counterexample \cite{shitov_counterexample_2018} lay outside this scope and thus might have still been a plausible counterexample.

[GO BACK TO QUESTION ON SECANT VARIETIES] Bernardi et al. \cite{bernardi_hitchhiker_2018} in 2018 comprehensively covered the literature and algorithms required for the application of secant varieties to questions about (generic) tensor decomposition, with a detailed explanation of the case of symmetric tensor decomposition and its relationship to the AH theorem.

[PROPOSITION 1. MAYBE THEOREM 1]Casarotti et al. \cite{casarotti_comons_2018} in 2018 proposed to improve the bounds for Comon's conjecture when working with flattenings, provided that equations for the secant varieties associated to the Veronese variety could be generated. Their result was for a tensor $T$ having rank $r < \binom{n + \lfloor \frac{d}{2} \rfloor}{n}$, the $\srank(T) = r$, in cases where such generating determinantal equations were known.

[THEOREM 1.3 - MAY NEED TO COVER ALL SECTIONS?]Finally in a paper from 2020, Seigal \cite{seigal_ranks_2020} demonstrated that for all "cubic surfaces", Comon's conjecture does hold. Here cubic surfaces were used to refer to the fact that under the identification with homogenous polynomials, order-$3$ symmetric tensors correspond to cubic polynomials, and thus their zero sets could be called cubic surfaces. The author also used flattenings of tensors (see \Cref{subsection:tensorflattening}), dividing their proof into cases depending on the number of singularities on the cubic surfaces in question.
\section{Attempts to Disprove the Conjecture}
% Strassen's conjecture: Determining the algorithmic complexity of matrix multiplication. Does there exist a tensor decomposition that computes two distinct matrix multiplications that costs less (has less rank) than the sum of the matrix multiplications individually? Additivity holds for bilinear maps \cite{landsberg_tensors_2012}
\section{Overview of Valid Cases of the Conjecture}\label{section:validoverview}
The current overview of cases where Comon's Conjecture is known to hold is as follows: \begin{itemize}
	\item For $T = \sum_{i=1}^{s} y_i^{\otimes d}$ and $d$ \textit{large enough} \cite{comon_symmetric_2008}. The large enough argument in the paper hinges on a classical result from algebraic geometry, that $d$ points in $\complex^n$ form the solution set of of polynomial equations of degree $\leq d$.
	\item For $T \in S^d(\complex^n)$ and $\srank(T) = 1, 2$ \cite{comon_symmetric_2008}.
	\item For $T \in S^d{\field^n}$, and $\rank(T) \leq \frac{dn - d + 1}{2}$ where $\field = \reals, \complex$ \cite{oeding_report_2008}.
	\item For $T \in S^d(\complex^n)$ for $(d,n) \in \{(3,2), (4,2), (3,3)\}$ \cite{friedland_remarks_2016}.
	\item In the case $d \geq 3,  \text{char}(\mathbb{F})  \geq 3, T \in S^d(\mathbb{F}^n)$ and $\rank(T) = \rank(T^{(i)})$ or $\rank(T^{(i)}) + 1$ (Here $T^{(i)}$ denotes the flattening (see \Cref{subsection:tensorflattening}) of $T$ along any mode. Since $T$ is symmetric all directions are equivalent.) \cite{friedland_remarks_2016}
	\item When $T \in S^3(\complex^n)$ and $\rank(T) \leq 5$ or $\srank(T) \leq 6$ \cite{friedland_remarks_2016}.
	\item For $T \in S^d(\mathbb{F}^n)$ where $\rank(T) \leq d$ where $\mathbb{F} = \reals , \complex$\cite{zhang_comons_2016}.
	\item If $T$ is a tensor lying on a tangential variety to a Veronese variety \cite{ballico_tensor_2013}.
	\item If $T \in \complex^2 \otimes \complex^n \otimes \complex^n$ \cite{buczynski_ranks_2013}.
	\item For $T$ being a Coppersmith-Winograd tensor \cite{landsberg_abelian_2017}.
	\item Whenever a tensor $T$ is cubic surface (real or complex) i.e. $T$ can be represented by a homogenous cubic polynomial $ \sum_{|\alpha|=3} \lambda_{\alpha}\mathbf{x}^{\alpha}$ \cite{seigal_ranks_2020}.
\end{itemize}


\end{document}
